% Formal checks of declared reconstructions; no proof of historical classification.
\section{Johnson's assertion basis and comparison cases}
\label{sec:p13-johnson-annex}

\begin{annexitem}{annex:p13.johnson.immediate}{thm:p13-johnson-la-immediate}{Immediate assertoric deductions}
\begin{proof}
Every propositional tautology with categorical formulas substituted for its
propositional letters is available by A0. Consequently modus ponens
permits antecedent interchange, contraposition, composition of implications,
and construction of a biconditional from its two directions. Each such
use below can be expanded by inserting the corresponding A0 conditional
and applying modus ponens; no categorical premise is being added.

A4 is \(Abc\to(Iba\to Iac)\). Substitute \(c=b\) to obtain
\(Abb\to(Iba\to Iab)\). A1 gives \(Abb\), so modus ponens gives
\(Iba\to Iab\). Interchange the letters \(a,b\) uniformly to obtain
\(Iab\to Iba\). Contraposition of \(Iba\to Iab\) yields
\(\neg Iab\to\neg Iba\). Replacing both expressions by their
defined \(E\)-forms yields \(Eab\to Eba\).

For subalternation substitute \(b=a\) into A4, giving
\(Aac\to(Iaa\to Iac)\). The propositional tautology that interchanges
the two antecedents gives \(Iaa\to(Aac\to Iac)\).
A2 and modus ponens give \(Aac\to Iac\). Substitute \(c=b\):
\(Aab\to Iab\). Contraposition gives
\(\neg Iab\to\neg Aab\), hence \(Eab\to Oab\).

Composition of \(Aab\to Iab\) with \(Iab\to Iba\) proves
\(Aab\to Iba\). Compose \(Eab\to Eba\) with the uniformly
substituted subalternation \(Eba\to Oba\) to obtain \(Eab\to Oba\).
This proves all six implications.

The first and third stated oppositions follow by definition:
\(Oab\) is \(\neg Aab\), and \(Eab\) is \(\neg Iab\).
For the other two use the A0 biconditional
\(\neg\neg p\leftrightarrow p\) and replace the definitional
expressions. Thus \(\neg Eab=\neg\neg Iab\) is equivalent to
\(Iab\), and \(\neg Oab=\neg\neg Aab\) is equivalent to \(Aab\).
Both directions of each equivalence are available by A0 and
definition interchange. All four biconditionals follow.
\end{proof}
\end{annexitem}

\begin{annexitem}{annex:p13.johnson.contrasted}{thm:p13-johnson-contrasted}{Two finite contrasted instances}
\begin{proof}
Use domain \(D=\{u,v\}\), \(u\ne v\). Interpret \(Axy\) as
inclusion of the nonempty extension of \(x\) in that of \(y\),
\(Ixy\) as nonempty intersection, \(Exy\) as empty intersection,
and \(Oxy\) as a nonempty difference.

In the first model put \(a=\{u,v\}, b=\{u\}, c=\{v\}\).
Then \(b\subseteq a\) and \(c\cap b=\varnothing\), so
\(Aba,Ecb\) hold. But \(c\subseteq a\), so \(Eca\) is false
and \(Oca\) is false. The model therefore defeats both negative
conclusions, not just a universal negative one.

In the second model put \(a=b=\{u\}\) and \(c=\{v\}\).
Again \(Aba,Ecb\) hold. Since \(c\cap a=\varnothing\) and
\(c\ne\varnothing\), \(Ica\) and \(Aca\) are false.
It defeats both positive conclusions. Each mentioned term is nonempty
in both models. A semantic consequence would have to hold in every
model of the premises; these two models defeat all four candidates.
\end{proof}
\end{annexitem}

\begin{annexitem}{annex:p13.johnson.lm}{thm:p13-johnson-lm-laws}{All four matrix laws}
\begin{proof}
Represent the four values bijectively by pairs of classical bits:
\(1=(1,1),2=(1,0),3=(0,1),4=(0,0)\).
The displayed implication table is coordinatewise classical implication.
Indeed antecedent \(1\) returns the consequent unchanged; antecedent
\(2\) gives row \(1,1,3,3\); antecedent \(3\) gives row
\(1,2,1,2\); and antecedent \(4\) gives row \(1,1,1,1\).
These are exactly the four displayed rows under the bijection.
The displayed negation complements both bits. The modal columns give
\(M(s,t)=(s,1)\) and \(L(s,t)=(s,0)\) for each of the four pairs.

Let \(p=(s,t),q=(u,v)\), and write
\(r=s\Rightarrow u\), \(z=t\Rightarrow v\), where
\(\Rightarrow\) is the classical bit implication. Then
\(p\to q=(r,z)\), while both \(Mp\to Mq\) and
\(Lp\to Lq\) equal \((r,1)\):
the second coordinate is respectively \(1\Rightarrow1\) or
\(0\Rightarrow0\), both \(1\).

It follows that
\((p\to q)\to(Mp\to Mq)\) and
\((p\to q)\to(Lp\to Lq)\) both have value
\((r\Rightarrow r,z\Rightarrow1)=(1,1)\).
Also \(L(p\to q)=(r,0)\). Its implication to either modal
consequent has value
\((r\Rightarrow r,0\Rightarrow1)=(1,1)\).
The designated pair is exactly value \(1\), so every assignment gives
the required value to all four formulas.

For the two propagation formulas, let the value of
\(p\to(q\to r)\) be \((f,g)\). Inserting \(L\) in either of the
stated positions leaves the first coordinate unchanged, since \(L\)
preserves the first bit. The second coordinate of
\(p\to(Lq\to Lr)\) is \(p_2\Rightarrow(0\Rightarrow0)=1\).
That of \(Lp\to(q\to Lr)\) is
\(0\Rightarrow(q_2\Rightarrow0)=1\).
Each consequent therefore has value \((f,1)\), and each outer
implication has value \((f\Rightarrow f,g\Rightarrow1)=(1,1)\).
\end{proof}
\end{annexitem}

\begin{annexitem}{annex:p13.johnson.modal}{thm:p13-johnson-modal-countercases}{Witness preservation and an S5 countermodel}
\begin{proof}
For the first claim reason at the world where both premises hold.
The existential premise supplies a witness \(d\) with \(Bd\) and
\(\Box\neg Cd\). Instantiate the universal premise at that same
\(d\) to obtain \(Bd\to Ad\). Modus ponens gives \(Ad\).
Thus \(Ad\land\Box\neg Cd\), and existential generalization gives
\(\exists x(Ax\land\Box\neg Cx)\). Equivalently a syntactic
existential-elimination proof uses a fresh parameter \(d\) not free in
the universal premise or final existential conclusion and discharges
the witness assumption. No inference moving a box across an actual
implication is used.

For the second claim let \(W=\{w_1,w_2\}\), let accessibility be
\(W\times W\), and use constant domain \(D=\{m,n\}\), \(m\ne n\).
The full relation is reflexive, symmetric, and transitive, so it is
an S5 frame. At both worlds \(A\) has extension \(\{m\}\) and
\(B\) has extension \(\{n\}\).
At \(w_1\), \(C=\{n\}\); at \(w_2\), \(C=\{m\}\).

At \(w_1\), the universal premise holds. Its instance for \(m\)
has false antecedent \(Cm\); its instance for \(n\) has true
antecedent \(Cn\) and true consequent \(\Box Bn\), since \(Bn\)
holds at both accessible worlds. The second premise has witness \(m\):
\(Am\) holds at both worlds and \(Bm\) at neither, giving
\(\Box Am\land\Box\neg Bm\).

The conclusion has no witness at \(w_1\). For \(m\),
\(\Box\neg Cm\) fails because \(Cm\) holds at \(w_2\).
For \(n\), \(\Box An\) fails because \(An\) holds at neither
world. These exhaust the domain. Both premises are true and the
conclusion false, proving the required failure.

For the mixed Darii claim retain the two-world full-accessibility
frame and domain \(\{m,n\}\). At \(w_1\) put \(A=B=\{m\}\);
at \(w_2\) put \(A=B=\{n\}\). Let \(C=\{m,n\}\) at both
worlds. At \(w_1\), each instance \(Bx\to\Box Cx\) holds because
\(\Box Cx\) holds for both individuals. The actual particular
premise has witness \(m\). But no individual is \(A\) in both
worlds, so neither can satisfy \(\Box(Ax\land Cx)\).
The existential conclusion is false. All three simple predicate
extensions are nonempty at both worlds.
\end{proof}
\end{annexitem}

\begin{annexitem}{annex:p13.johnson.geach}{thm:p13-johnson-geach-countercases}{Finite G-system interpretations}
\begin{proof}
Use ordinary nonempty-term extensional meanings for the assertoric
categoricals. For Darii LXL let \(D=\{u,v\}\), \(u\ne v\), and set
\[
 a=\{u,v\},\quad\lambda a=\{v\},\quad\mu a=\{u,v\},
 \qquad b=\lambda b=\mu b=\{u\},
 \qquad c=\lambda c=\mu c=\{u\}.
\]
All simple and complex terms are nonempty, and
\(\lambda x\subseteq x\subseteq\mu x\) holds for each.
The first premise \(LAbc\) translates to \(Ab\lambda c\),
true since \(\{u\}\subseteq\{u\}\).
The second \(Iab\) is true with witness \(u\).
The conclusion \(LIac\) translates to \(I\lambda a c\),
false since \(\{v\}\cap\{u\}=\varnothing\).

For Cesare LLL again use \(D=\{u,v\}\), and put
\[
 a=\lambda a=\mu a=\{u\},\qquad
 b=\lambda b=\mu b=\{u\},\qquad
 c=\lambda c=\{v\},\quad\mu c=\{u,v\}.
\]
The inclusion constraints and nonemptiness again hold.
The first premise \(LEcb\) is \(Ec\mu b\), true because
\(\{v\}\cap\{u\}=\varnothing\).
The second premise \(LAab\) is \(Aa\lambda b\), true by inclusion.
The conclusion \(LEac\) is \(Ea\mu c\), false since
\(\{u\}\cap\{u,v\}=\{u\}\).
Each is therefore a countermodel of the stipulated translation.
\end{proof}
\end{annexitem}
